\documentclass[11pt,a4paper]{article}

\usepackage[utf8]{inputenc}
\usepackage[T1]{fontenc}
\usepackage{geometry}
\geometry{margin=1in}
\usepackage{hyperref}
\usepackage{setspace}
\usepackage{parskip}

% Remove page numbers for the cover letter
\pagestyle{empty}

\begin{document}

\begin{flushright}
\textbf{Connor Forbes}\\
Institute for Evidence-Based Healthcare\\
Bond University\\
Gold Coast, Queensland, Australia\\
Email: cforbes@bond.edu.au\\
\today
\end{flushright}

\vspace{2em}

\begin{flushleft}
\textbf{Prof. Dimitris Mavridis and Prof. Terri D. Pigott}\\
Editors-in-Chief\\
\textit{Research Synthesis Methods}\\
Cambridge University Press
\end{flushleft}

\vspace{2em}

Dear Prof. Mavridis and Prof. Pigott,

\vspace{1em}

\textbf{Re: Submission of manuscript ``MechaScreener: Large Language Model-Based Automated Screening for Systematic Reviews and Research''}

\vspace{1em}

Please find enclosed our manuscript, which we are submitting for consideration as an Original Research Article in \textit{Research Synthesis Methods}.

Systematic reviews remain the gold standard of evidence synthesis, yet the manual screening of titles and abstracts creates a severe, resource-intensive bottleneck. While previous automated tools and Large Language Models (LLMs) have shown significant promise, they have historically struggled to consistently achieve the near-perfect recall required to ensure no critical studies are missed. Our manuscript addresses this pivotal methodological challenge.

In this study, we introduce and evaluate MechaScreener, a zero-shot LLM-based automated screening tool that requires no prior training data. By applying user-provided PICOT or general eligibility criteria, the tool assigns deterministic inclusion probability scores to each reference. We robustly evaluated the tool across 10 diverse Cochrane systematic review datasets (encompassing both randomised controlled trials and non-RCTs) containing over 58,000 references. Crucially, MechaScreener achieved a perfect mean recall (1.00 or 100\%) while safely eliminating an average of 61\% of irrelevant literature during the initial screening phase.

We believe this paper is an excellent fit for the readership of \textit{Research Synthesis Methods}. It presents a validated methodological advancement in evidence synthesis automation. By demonstrating how an LLM can be conservatively tuned to function as a highly reliable ``first-pass'' filter, our findings offer reviewers a practical mechanism to substantially reduce manual workloads without compromising the rigorous safety standards of systematic reviews.

We confirm that this manuscript is original, has not been published elsewhere, and is not currently under consideration by another journal. All authors have reviewed and approved the final manuscript. In the spirit of methodological transparency and open science, the dataset, source code, and scripts used for our analysis have been made publicly available on GitHub. As declared in the manuscript, the authors are responsible for the development of the MechaScreener tool and warmly welcome independent testing to verify our findings.

Thank you for your time and consideration. We look forward to the opportunity to receive feedback from your reviewers.

\vspace{2em}

Sincerely,

\vspace{3em}

\textbf{Connor Forbes}\\
Corresponding Author\\
On behalf of co-authors: Matt Carter, Carly Hudson, and Justin Clark.

\end{document}